%===========================================================
% FINAL MANUSCRIPT (RevTeX4-2) -- with FBS^3R bridge section inserted
% Title: Quasi-local diffeomorphism charges in a closed S^3 universe
%===========================================================
\documentclass[
  reprint,
  amsmath,amssymb,
  aps,
  prd,
  superscriptaddress,
  nofootinbib,
  longbibliography
]{revtex4-2}
%------------------ Packages ------------------
\usepackage[T1]{fontenc}
\usepackage[utf8]{inputenc}
\usepackage{lmodern}
\usepackage{microtype}
\usepackage{bm}
\usepackage{mathtools}
\usepackage{amsthm}
\usepackage{siunitx}
\sisetup{per-mode = symbol}
\usepackage{graphicx}
\usepackage{booktabs}
% --- Plot generated inside LaTeX (no external figure file) ---
\usepackage{tikz}
\usepackage{pgfplots}
% Publisher fallback: if pgfplots is disabled, replace Fig.~1 with an external PDF figure.
\pgfplotsset{compat=1.18}
\usepackage[hidelinks]{hyperref}
\usepackage{orcidlink}
\makeatletter
\AtBeginDocument{%
  \@ifpackageloaded{hyperref}{\hypersetup{hidelinks}}{}%
}
\makeatother
%------------------ Macros ------------------
\newcommand{\dd}{\mathrm{d}}
\newcommand{\phifield}{\varphi} % glyph (phi), if needed
\newcommand{\goldconst}{\ensuremath{\phi_{\!\mathrm{g}}}} % golden ratio constant symbol
\newcommand{\gold}{\goldconst}
\newcommand{\phig}{\goldconst} % shorthand for golden ratio constant
\newcommand{\Sthree}{\ensuremath{S^3}}
\newcommand{\M}{\mathcal{M}}
\newcommand{\Sig}{\Sigma}
\newcommand{\del}{\partial}
\newcommand{\Lie}{\mathcal{L}}
\newcommand{\on}{\approx}
\newcommand{\G}{G}
\newcommand{\Lform}{\bm{L}}
\newcommand{\tht}{\bm{\theta}}
\newcommand{\om}{\bm{\omega}}
\newcommand{\J}{\bm{J}}
\newcommand{\Q}{\bm{Q}}
\newcommand{\Bnd}{\bm{B}}
\newcommand{\eps}{\bm{\epsilon}}
%------------------ Theorem environments ------------------
\newtheorem{theorem}{Theorem}
\newtheorem{prop}{Proposition}
\newtheorem{claim}{Claim}
\newtheorem{cor}{Corollary}
\begin{document}
% Use \texorpdfstring to avoid hyperref PDF-string warnings due to math in title
\title{\texorpdfstring{Quasi-local diffeomorphism charges in a closed $\Sthree$ universe:\\
explicit Iyer--Wald charges, the Gibbons--Hawking--York term,\\
and $\goldconst$-level boundary scanning}{Quasi-local diffeomorphism charges in a closed S^3 universe: explicit Iyer--Wald charges, the Gibbons--Hawking--York term, and golden-ratio level boundary scanning}}
\author{Andrei Preece\orcidlink{0009-0004-5978-795X}}
\affiliation{Independent Researcher, Berlin, Germany}
\author{Boris Batenin\orcidlink{0009-0000-7027-8583}}
\affiliation{Independent Researcher, Berlin, Germany}
\date{\today}
\begin{abstract}
Diffeomorphism invariance implies Noether identities and first-class constraints rather than ordinary global conserved charges.
In a compact closed universe with Cauchy slice $\Sigma \simeq S^{3}$ one has $\partial\Sigma=\varnothing$, hence the global Hamiltonian generator associated with gauge diffeomorphisms vanishes.
Nevertheless, physically meaningful generators arise quasi-locally for subregions $D\subset\Sigma$ with $\partial D\neq\varnothing$ and are entirely encoded by boundary integrals.
We provide a self-contained covariant phase space derivation in general relativity with an explicit Iyer--Wald Noether charge two-form $Q_\xi$ and an explicit boundary three-form $B$ corresponding to the Gibbons--Hawking--York term.
For closed FRW geometry we compute the Brown--York energy of geodesic balls in $\Sigma\simeq S^{3}$ and obtain a closed expression in terms of the areal radius $R$ and curvature radius $a(t)$.
We relate this boundary energy to the FRW Misner--Sharp mass and clarify that the two coincide only after additional dynamical input (e.g.\ Friedmann constraints and/or a choice of evolution vector field), since Brown--York is tied to a chosen slice and evolution vector field, whereas Misner--Sharp is the invariant quasi-local mass in spherical symmetry and is often associated with the Kodama flow.\cite{Kodama1980,Hayward1996}
We also highlight constant-mean-curvature slicing and York time as a canonical choice of relational clock in compact universes.
Finally, we introduce a $\goldconst$-discretized boundary scan (a golden-ratio discretization of radii, not a dynamical scalar field) $R_n = R_0\goldconst^{n}$ as a structured quasi-local bookkeeping device; a model-dependent $\zeta_{\goldconst}$ spectral overlay and a modular reflection map $n\leftrightarrow n'$ are recorded separately in an appendix.
\end{abstract}
\maketitle
\section{Introduction}
In diffeomorphism-invariant theories, conservation laws take a different structural form from those of global symmetries.
Noether's second theorem implies identities among the Euler--Lagrange equations and yields first-class constraints.
In general relativity (GR), the full compact-universe logic may be summarized in the one-line chain
\begin{equation}
\boxed{
\begin{gathered}
\mathrm{Diff}(\M)\Rightarrow \text{Noether II}\Rightarrow \J_{\xi}\approx \dd \Q_{\xi}\\
H_{\xi}=\int_{\partial\Sigma}\bigl(\Q_{\xi}-\xi\cdot\Bnd\bigr)
\Rightarrow \partial\Sigma=\varnothing\\
\Rightarrow H_{\xi}[\Sigma]=0 (\text{gauge }\xi)
\end{gathered}
}
\label{eq:masterChainBox}
\end{equation}
where $\approx$ denotes equality on-shell.
The key point for cosmology is that a spatially closed universe with compact Cauchy slice $\Sigma\simeq S^{3}$ has no boundary.
Hence there is no nontrivial \emph{global} boundary generator for \emph{gauge} diffeomorphisms.
Physical information is encoded (i) in constraints, and (ii) in quasi-local boundary charges for subregions $D\subset\Sigma$.
This paper makes these statements explicit in the covariant phase space formalism and computes the resulting quasi-local energy in closed FRW geometry.
Throughout, $\xi^a$ denotes a smooth vector field used as an \emph{evolution/gauge parameter} in the covariant phase space formalism.
Unless $\xi$ is a genuine symmetry (e.g.\ Killing), the associated quasi-local generator need not be conserved under time evolution; the construction yields a meaningful Hamiltonian generator on a chosen phase space with specified boundary conditions, not automatically a globally conserved charge in the Noether-I sense.
\section{From Noether I to Noether II, boundary charges, and ADM}
It is useful to situate the chain \eqref{eq:masterChainBox} within the more familiar Noether-I picture.
For ordinary global symmetries (finite-dimensional parameters), Noether's first theorem yields a conserved current
$\nabla_\mu j^\mu = 0$ and an associated charge obtained by integrating $j^\mu$ over a Cauchy slice.
In that setting the charge is genuinely ``bulk'': its conservation follows from the local continuity equation together with Stokes' theorem.
Diffeomorphisms, by contrast, are a local gauge redundancy with an arbitrary spacetime-dependent parameter $\xi^\mu(x)$.
Noether's second theorem therefore applies: invariance under $\xi^\mu(x)$ produces identities among the Euler--Lagrange equations
(Bianchi-type relations) and, in the Hamiltonian language, first-class constraints.
In covariant phase space (CPS) this shows up as a strengthening of ``current conservation'': the diffeomorphism Noether current is on-shell exact,
\begin{equation}
  \J_\xi \on \dd \Q_\xi ,
\end{equation}
so the would-be charge is necessarily encoded by a surface integral.
More precisely, writing the Lagrangian as a top form $\Lform$ and its variation as
\begin{equation}
\delta\Lform = E(\phi)\,\delta\phi + \dd\tht(\phi;\delta\phi),
\end{equation}
one defines
\begin{equation}
\J_\xi := \tht(\phi;\Lie_\xi\phi)-\xi\cdot\Lform.
\end{equation}
The CPS symplectic current $\om=\delta\tht$ then gives the variation of the Hamiltonian generator as a boundary term,
\begin{equation}
  \delta H_\xi[D]=\int_{\partial D}\Bigl(\delta\Q_\xi-\xi\cdot\tht(\phi;\delta\phi)\Bigr),
\label{eq:dHxi}
\end{equation}
for any spatial region $D\subset\Sigma$.
If the boundary conditions and the choice of phase space imply integrability,
\begin{equation}
\xi\cdot\tht(\phi;\delta\phi)=\delta(\xi\cdot\Bnd)
\quad \text{on the chosen phase space,}
\label{eq:integrability}
\end{equation}
then one obtains an actual generator
\begin{equation}
  H_\xi[D]=\int_{\partial D}\bigl(\Q_\xi-\xi\cdot\Bnd\bigr),
\label{eq:Hxi}
\end{equation}
up to an additive constant.
In general, non-integrability can occur (physically: boundary flux); addressing it may require additional structure (e.g.\ Wald--Zoupas-type prescriptions\cite{WaldZoupas2000}). In this work we restrict attention to boundary data for which \eqref{eq:integrability} holds. Concretely, we assume Dirichlet boundary conditions for the induced metric on the timelike worldtube swept out by $\partial D$, so that the GHY choice makes the Hamiltonian integrable on the chosen phase space.
The ADM formulation makes the same statement in canonical language.
Decomposing $\xi^\mu=N n^\mu+N^\mu$ into lapse $N$ and shift $N^i$, the canonical generator takes the form
\begin{equation}
  H_{\rm ADM}[N,N^i]
  =\int_{\Sigma}\bigl(N\,\mathcal{H}+N^i\mathcal{H}_i\bigr)\;+\;\oint_{\partial\Sigma}\mathcal{B}[N,N^i].
\end{equation}
Here $\mathcal{H}\approx0$ and $\mathcal{H}_i\approx0$ are the first-class Hamiltonian and momentum constraints implied by Noether II,
while $\mathcal{B}$ is the boundary functional needed for differentiability and for nonzero charges when a boundary exists.
On-shell the bulk term vanishes and the generator is purely boundary-supported. Therefore, in a closed universe with compact $\Sigma\simeq S^{3}$ (or more generally any compact $\Sigma$ without boundary),
\begin{equation}
  \partial\Sigma=\varnothing
  \qquad\Longrightarrow\qquad
  H_\xi[\Sigma]=0
\end{equation}
for all smooth \emph{gauge} diffeomorphisms, while quasi-local charges can still be defined for subregions $D\subset\Sigma$ with $\partial D\neq\varnothing$.
\subsection{Closed FRW in one line: \texorpdfstring{$\partial\Sigma=\varnothing\Rightarrow H_\xi[\Sigma]=0\ \text{(gauge }\xi\text{)}$}{dSigma=empty => H_xi[Sigma]=0}}
In a closed FRW cosmology $\M=\mathbb{R}\times S^{3}$ one may write
\begin{equation}
\dd s^2=-\dd t^2+a(t)^2\,\dd\Omega_3^2,
\qquad
\Sigma_t\simeq S^3,
\end{equation}
so every Cauchy slice is compact and boundaryless: $\partial\Sigma_t=\varnothing$. Therefore
\begin{equation}
H_\xi[\Sigma_t] = \int_{\partial\Sigma_t}\bigl(\Q_\xi-\xi\cdot\Bnd\bigr)=0,
\end{equation}
for any smooth diffeomorphism parameter $\xi$ interpreted as a gauge direction on the phase space.
\subsection{Explicit \texorpdfstring{$\Q_\xi$}{Q_xi} (Komar/Iyer--Wald) and a quasi-local \texorpdfstring{$S^2\subset S^3$}{S2 subset S3} example}
For Einstein--Hilbert gravity, the Iyer--Wald Noether charge two-form is
\begin{equation}
(\Q_\xi)_{ab}=-\frac{1}{16\pi G}\,\epsilon_{abcd}\,\nabla^c\xi^d.
\label{eq:QEH}
\end{equation}
Even though the global slice has no boundary, quasi-local charges are nontrivial for subregions.
Take a geodesic ball $D\subset\Sigma_t\simeq S^{3}$ with boundary $\partial D\simeq S^{2}$ (e.g.\ $\chi=\chi_0$ in hyperspherical coordinates), and choose $\xi=u=\partial_t$ (the comoving time flow).
Then the quasi-local generator is a boundary integral on $\partial D$,
\begin{equation}
H_u[D] = \int_{\partial D}\bigl(\Q_u-u\cdot \Bnd\bigr).
\end{equation}
With a suitable boundary term $\Bnd$ and a reference subtraction fixing the zero-point, this reproduces the Brown--York energy computed below.
Since $u$ is not generally Killing in FRW, $H_u[D]$ should be viewed as a quasi-local Hamiltonian associated with the chosen slicing/evolution, not as an a priori conserved charge.
\section{Covariant phase space preliminaries}
We follow the covariant phase space formalism of Lee--Wald and the diffeomorphism-charge construction of Iyer--Wald.\cite{LeeWald1990,IyerWald1994}
Let $\Lform(\phi)$ be a Lagrangian $4$-form on spacetime $\M$.
Its variation is
\begin{equation}
\delta\Lform = E(\phi)\,\delta\phi + \dd\tht(\phi;\delta\phi),
\end{equation}
where $E(\phi)=0$ are equations of motion and $\tht$ is the symplectic potential $3$-form.
The symplectic current $3$-form is
\begin{equation}
\om(\phi;\delta_1\phi,\delta_2\phi)=\delta_1\tht(\phi;\delta_2\phi)-\delta_2\tht(\phi;\delta_1\phi).
\end{equation}
For a vector field $\xi^{a}$, diffeomorphism invariance implies $\delta_\xi\Lform=\Lie_\xi\Lform=\dd(\xi\cdot\Lform)$.
\section{Noether current, Noether II, and boundary Hamiltonians}
Define the Noether current $3$-form
\begin{equation}
\J_\xi := \tht(\phi;\Lie_\xi\phi)-\xi\cdot\Lform.
\end{equation}
A standard identity yields $\dd\J_\xi=-E(\phi)\Lie_\xi\phi$, so $\dd\J_\xi\on 0$.
For gauge symmetries such as diffeomorphisms, Noether II strengthens this to on-shell exactness:
\begin{equation}
\J_\xi \on \dd\Q_\xi,
\end{equation}
for some locally constructed $2$-form $\Q_\xi$.
In the covariant phase space formalism, the variation of the Hamiltonian generator associated with $\xi$ over a spatial region $D\subset\Sigma$ reduces to a boundary expression \eqref{eq:dHxi}.
When \eqref{eq:integrability} holds for the chosen boundary conditions, the integrable generator is given by \eqref{eq:Hxi}.
\section{Explicit GR formulas: Iyer--Wald \texorpdfstring{$\Q_\xi$}{Q_xi} and the GHY boundary form \texorpdfstring{$\Bnd$}{B}}
We now specialize to GR (with optional $\Lambda$). Let $\eps$ denote the spacetime volume $4$-form. The Einstein--Hilbert Lagrangian is
\begin{equation}
\Lform_{\rm EH}=\frac{1}{16\pi G}(R-2\Lambda)\,\eps.
\end{equation}
In Iyer--Wald conventions, the symplectic potential may be written as
\begin{equation}
\tht_{abc}(g;\delta g)=\frac{1}{16\pi G}\,\epsilon_{dabc}\Bigl(\nabla_e\delta g^{de}-\nabla^d\delta g\Bigr),
\end{equation}
where $\delta g := g_{mn}\delta g^{mn}$ and indices are moved with $g_{ab}$.
The Iyer--Wald Noether charge $2$-form is \eqref{eq:QEH}.
To make the variational principle well-posed on manifolds with a spacetime boundary hypersurface $\mathcal{B}$ under Dirichlet boundary conditions (fixed induced metric on $\mathcal{B}$), one adds the Gibbons--Hawking--York term,\cite{GibbonsHawking1977,York1972}
\begin{equation}
S_{\rm GHY}=\frac{1}{8\pi G}\int_{\mathcal B}\sqrt{|h|}\,K\,\dd^3x,
\end{equation}
corresponding to the boundary $3$-form
\begin{equation}
\Bnd_{\rm GHY}=\frac{1}{8\pi G}\,K\,\eps_{\mathcal B}.
\end{equation}
Here $K$ is the trace of the extrinsic curvature of the \emph{3D} boundary $\mathcal{B}\subset\M$ in spacetime, and $\eps_{\mathcal B}$ is its induced volume form.
In quasi-local Hamiltonians one ultimately evaluates boundary functionals on \emph{2-surfaces} $\partial D\subset\Sigma$; a reference subtraction (or counterterm) is additionally required to fix the energy zero-point, as in the Brown--York construction below.
\section{Closed FRW geometry and quasi-local energy on \texorpdfstring{$\Sigma\simeq S^3$}{Sigma~S3}}
Consider a closed FRW slicing with spatial metric on $\Sigma_t\simeq S^3$:
\begin{equation}
\dd\ell^2=a(t)^2\left(\dd\chi^2+\sin^2\chi\,\dd\Omega^2\right),\qquad \chi\in[0,\pi].
\end{equation}
Let $D(\chi_0)$ be a geodesic ball $0\le\chi\le\chi_0$ with boundary $\partial D\simeq S^2$ at $\chi=\chi_0$. The areal radius and area are
\begin{equation}
R(\chi_0)=a(t)\sin\chi_0,\qquad A(\chi_0)=4\pi R^2,
\end{equation}
and the volume is
\begin{equation}
V(\chi_0)=2\pi a(t)^3\bigl(\chi_0-\sin\chi_0\cos\chi_0\bigr).
\end{equation}
The Brown--York energy of $D$ is
\begin{equation}
E_{\rm BY}=\frac{1}{8\pi G}\int_{\partial D}(k_0-k)\,\dd S,
\label{eq:EBYdef}
\end{equation}
where $k$ is the trace of the extrinsic curvature of $\partial D$ embedded in $\Sigma$ (with outward unit normal in $\Sigma$) and $k_0$ is a reference term fixing the zero-point.
A standard choice is the Euclidean reference: $k_0$ is taken as the mean curvature of an isometric embedding of the same intrinsic 2-metric into $\mathbb{R}^3$, giving $k_0=2/R$. For the round $S^2$ boundary with intrinsic metric $R^2 d\Omega^2$, the standard isometric embedding into $\mathbb{R}^3$ exists and yields this value.
For $\chi=\chi_0$ spheres in $\Sigma\simeq S^3$,
\begin{equation}
k(\chi_0)=\frac{2}{a}\cot\chi_0=\frac{2}{R}\sqrt{1-\frac{R^2}{a^2}},\qquad
k_0=\frac{2}{R}.
\end{equation}
Substituting into \eqref{eq:EBYdef} yields the closed expression
\begin{equation}
E_{\rm BY}(R,a)=\frac{R}{G}\left(1-\sqrt{1-\frac{R^2}{a^2}}\right),
\qquad 0\le R\le a(t).
\label{eq:EBYclosed}
\end{equation}
This result was previously obtained in the equivalent coordinate system by Afshar \cite{Afshar2009}.
\begin{figure}[htbp]
\centering
\begin{tikzpicture}
\begin{axis}[
  width=0.85\linewidth,
  height=0.50\linewidth,
  xlabel={$r=R/a$},
  ylabel={$G E_{\rm BY}/a$},
  xmin=0, xmax=0.7,
  ymin=0,
  domain=0:0.7,
  samples=300,
  axis lines=left,
  legend style={draw=none, at={(0.03,0.97)}, anchor=north west},
]
\addplot+[mark=none] { x*(1 - sqrt(1 - x^2)) };
\addlegendentry{Exact}
\addplot+[mark=none, dashed] { x^3/2 };
\addlegendentry{$r^3/2$}
\end{axis}
\end{tikzpicture}
\caption{Brown--York quasi-local energy in closed FRW cosmology:
$G E_{\rm BY}/a$ versus $r=R/a$ (solid curve).
The small-ball approximation $r^3/2$ is shown dashed for comparison.}
\label{fig:EBY}
\end{figure}
At the equator ($R\to a$) one has $E_{\rm BY}\to R/G$.
For small balls $R\ll a$,
\begin{equation}
E_{\rm BY}(R,a)=\frac{R^3}{2Ga^2}+O\!\left(\frac{R^5}{a^4}\right).
\end{equation}
This small-ball limit highlights that, even before invoking dynamics, the positive spatial curvature of $S^3$ induces a nontrivial quasi-local energy relative to the Euclidean reference.
\section{Misner--Sharp mass, Friedmann constraint, and comparison with Brown--York}
In spherically symmetric spacetimes, the Misner--Sharp mass provides an invariant quasi-local energy:
It is often associated with the Kodama vector field, which supplies a preferred energy current in spherical symmetry.
\begin{equation}
M_{\rm MS}(R)=\frac{R}{2G}\left(1-g^{ab}\nabla_aR\,\nabla_bR\right).
\end{equation}
For FRW, this becomes
\begin{equation}
M_{\rm MS}(R)=\frac{R^3}{2G}\left(H^2+\frac{k}{a^2}\right),
\qquad H:=\frac{\dot a}{a}.
\label{eq:MMSFRW}
\end{equation}
For the closed case ($k=+1$),
\begin{equation}
M_{\rm MS}(R)=\frac{R^3}{2G}\left(H^2+\frac{1}{a^2}\right).
\end{equation}
The Friedmann constraint reads
\begin{equation}
H^2+\frac{k}{a^2}=\frac{8\pi G}{3}\rho+\frac{\Lambda}{3},
\end{equation}
hence
\begin{equation}
M_{\rm MS}(R)=\frac{4\pi}{3}R^3\rho+\frac{\Lambda R^3}{6G}.
\label{eq:MMSrho}
\end{equation}
Equation \eqref{eq:EBYclosed} depends only on the intrinsic/extrinsic geometry of the 2-surface $\partial D$ as embedded in the chosen spatial slice $\Sigma_t$ (plus the chosen reference subtraction), while \eqref{eq:MMSFRW} is a spacetime-covariant quantity depending on $\nabla_a R$ and thus on expansion via $H$.
Accordingly, $E_{\rm BY}$ and $M_{\rm MS}$ should be viewed as complementary quasi-local energies: $E_{\rm BY}$ is naturally tied to a slicing and boundary conditions, whereas $M_{\rm MS}$ is invariant in spherical symmetry and directly relates to matter content through \eqref{eq:MMSrho}.
A precise identification between them requires additional dynamical input and/or a preferred evolution vector field, and is not automatic at the level of purely kinematical definitions.
\section{CMC slicing, York time, and relational observables in a compact universe}
In a compact universe, $H_\xi[\Sigma]=0$ expresses that the total Hamiltonian is a sum of first-class constraints and vanishes on-shell.
To interpret dynamics, one introduces a relational clock.
For general discussions of relational clocks and relational observables in Hamiltonian constrained systems, see e.g.\ \cite{BrownKuchar1995,Dittrich2006}.
A widely used choice is constant-mean-curvature (CMC) slicing. Let $h_{ij}$ be the induced metric on $\Sigma_t$ and $K_{ij}$ its extrinsic curvature in spacetime; define the mean curvature
\begin{equation}
K := h^{ij}K_{ij}.
\end{equation}
CMC slicing imposes $K=\mathrm{const}$ on each $\Sigma_t$.
For standard FRW slicings, $K$ is spatially constant automatically by homogeneity and isotropy, so the usual cosmological foliation is a CMC foliation.
York time is then defined (up to conventions) by
\begin{equation}
T_Y := -\frac{1}{3}K.
\end{equation}
In compact universes, such a choice supplies a canonical relational time variable compatible with constraint dynamics, providing a clean way to parametrize evolution despite the vanishing of the global Hamiltonian generator.
\section{\texorpdfstring{$\goldconst$}{golden-ratio}-level boundary scanning and a one-block correspondence map}
Define a $\goldconst$-discretized family of areal radii (where $\goldconst := (1+\sqrt{5})/2$ is the golden ratio)
\begin{equation}
\begin{aligned}
R_n &= R_0\goldconst^{n},\\
n   &= \frac{\log(R/R_0)}{\log\goldconst}.
\end{aligned}
\end{equation}
restricted in $\Sigma\simeq S^3$ by $0\le R_n\le a(t)$.
This discretization is introduced purely as a structured quasi-local bookkeeping device and does not represent a claim that GR dynamically selects the constant $\goldconst$.
Then
\begin{equation}
\begin{aligned}
E_{\rm BY}(n) &= \frac{R_n}{G}\left(1-\sqrt{1-\frac{R_n^2}{a^2}}\right),\\
M_{\rm MS}(n) &= \frac{R_n^3}{2G}\left(H^2+\frac{1}{a^2}\right).
\end{aligned}
\end{equation}
%===========================================================
% INSERTED BRIDGE SECTION (FBS^3R core summary)
%===========================================================
\section{Discrete \texorpdfstring{$\Sthree$}{S3} evolution and golden-ratio consistency (summary of the FBS$^{3}$R core)}
\label{sec:FBS3Rbridge}
We briefly summarize the structural assumptions of the discrete evolution model developed in a companion discrete-framework note,
restricting to the minimal elements needed to interface with the quasi-local diffeomorphism-charge analysis of
Secs.~II--VI. The purpose of this section is \emph{not} to re-derive the FBS$^{3}$R construction, but only to record its
postulates and to make the compatibility with covariant phase space (CPS) and GR explicit.
\subsection{Postulates (minimal)}
\label{subsec:FBS3Rpostulates}
\paragraph{(P1) Topological ontology.}
At each discrete stage the spatial manifold is compact, simply connected, and boundaryless:
\begin{equation}
\Sigma_n \simeq S^3,
\qquad
\partial\Sigma_n=\varnothing,
\qquad
\Sigma_n \cong \Sigma_{n+1}.
\end{equation}
No topology change is assumed.
\paragraph{(P2) Causal continuity (internal evolution).}
Evolution is internal (no external embedding space is required) and is represented as a discrete update
\begin{equation}
\text{state}_n \rightarrow \text{state}_{n+1},
\end{equation}
with the spatial topology fixed as in (P1).
\paragraph{(P3) Fixed total energy.}
A conserved total energy is postulated,
\begin{equation}
E_{\text{total}}=\text{const},
\qquad
\rho_n=\frac{E_{\text{total}}}{V_n},
\end{equation}
where $V_n$ is the model-assigned volume associated with the radius parameter introduced below.
(When comparing to GR expressions on $\Sigma\simeq S^3$, one may instead use the geometric $S^3$ volume of a
geodesic ball, cf.\ Sec.~VII; the present subsection merely records the postulated bookkeeping used in a companion discrete-framework note.)
\subsection{Discrete geometric scaling}
\label{subsec:FBS3Rscaling}
The discrete radius is taken to scale as
\begin{equation}
R_n = \ell_P\,\goldconst^{N_n},
\qquad
\goldconst:=\frac{1+\sqrt{5}}{2},
\label{eq:FBS3R_Rn}
\end{equation}
with integer level $N_n$ (possibly time-dependent through the discrete evolution).
A corresponding volume scaling is assigned in the FBS$^{3}$R model as $V_n\propto R_n^3$, hence
\begin{equation}
\rho_n \propto \goldconst^{-3N_n}.
\end{equation}
This provides a compact encoding of density scaling compatible with the fixed-energy postulate (P3).
\subsection{Quantum surface growth and Fibonacci asymptotics}
\label{subsec:FBS3Rfibonacci}
The number of ``surface quanta'' obeys a Fibonacci recursion,
\begin{equation}
N_{n+1}=N_n+N_{n-1},
\end{equation}
so that asymptotically
\begin{equation}
\lim_{n\to\infty}\frac{N_{n+1}}{N_n}=\goldconst.
\end{equation}
In particular, this yields an asymptotic $\goldconst$-ratio between successive radii,
\begin{equation}
\frac{R_{n+1}}{R_n}
=\goldconst^{\,N_{n+1}-N_n}
\;\;\xrightarrow[n\to\infty]{}\;\;
\goldconst,
\label{eq:FBS3R_phiConsistency}
\end{equation}
establishing golden-ratio consistency between (i) discrete geometric scaling, (ii) Fibonacci growth, and (iii) surface-based
bookkeeping of density levels.
(a companion discrete-framework note also introduces the normalization $F_n=\goldconst^n/\sqrt{5}$; we do not need it for the GR/CPS
derivations in the present paper.)
\subsection{Compatibility with diffeomorphism invariance and CPS/GR charges}
\label{subsec:FBS3R_CPS_compat}
Because each $\Sigma_n$ is diffeomorphic to $S^3$ and is globally boundaryless, the Hamiltonian generators of
\emph{gauge} diffeomorphisms vanish on the full slice:
\begin{equation}
H_\xi[\Sigma_n]=\int_{\partial\Sigma_n}\bigl(\Q_\xi-\xi\cdot\Bnd\bigr)=0.
\end{equation}
This is precisely the compact-universe statement emphasized in Eq.~\eqref{eq:masterChainBox}.
Nevertheless, for subregions $D\subset\Sigma_n$ with $\partial D\neq\varnothing$, quasi-local Hamiltonians remain
well-defined and are encoded by boundary integrals,
\begin{equation}
H_\xi[D]=\int_{\partial D}\bigl(\Q_\xi-\xi\cdot\Bnd\bigr),
\end{equation}
with $(\Q_\xi,\Bnd)$ as constructed explicitly in GR in Secs.~V--VI.
Thus, the discrete $\goldconst$-driven scaling recorded in \eqref{eq:FBS3R_Rn}--\eqref{eq:FBS3R_phiConsistency}
does not conflict with diffeomorphism invariance: it specifies a relational/discrete parametrization of evolution
within a fixed diffeomorphism class $\Sigma\simeq S^3$, while the CPS/GR charge algebra and the statement
$H_\xi[\Sigma]=0$ remain intact.
\subsection{York time as a relational parameter}
\label{subsec:FBS3R_york}
In compact $\Sigma\simeq S^3$ settings, dynamics is naturally described in relational terms.
A standard canonical choice is constant-mean-curvature (CMC) slicing and York time,
\begin{equation}
T_Y := -\frac{1}{3}K,
\end{equation}
with $K$ the mean curvature of the embedding of $\Sigma$ in spacetime (Sec.~IX).
Within this viewpoint, the discrete index $n$ (or level $N_n$) may be treated as a monotonic function of $T_Y$,
\begin{equation}
n=n(T_Y),
\end{equation}
providing a relational clock compatible with the vanishing of global gauge generators on a closed slice.
%===========================================================
% END INSERT
%===========================================================
\section{Conclusions}
In a compact closed universe $\Sigma\simeq S^3$, diffeomorphism invariance implies that global gauge diffeomorphism generators vanish: $H_\xi[\Sigma]=0$.
Nontrivial physics is carried by constraints and by quasi-local boundary charges for subregions.
We made this explicit in GR by providing the closed-form Iyer--Wald charge $\Q_\xi$ and the GHY boundary form $\Bnd$, together with the standard integrability condition needed to define a Hamiltonian generator.
For closed FRW geometry, we computed the Brown--York energy of geodesic balls and related it to the Misner--Sharp mass and the Friedmann constraint, emphasizing the complementary (slice-adapted vs spacetime-covariant) character of these notions.
CMC slicing and York time supply a canonical relational clock compatible with compact-universe constraint dynamics.
Finally, a $\goldconst$-discretized boundary scan offers a structured quasi-local bookkeeping scheme; model-dependent spectral overlays can be layered without altering the GR derivations.
\appendix
\section{Derivation notes and ambiguities}
The covariant phase space formalism includes standard ambiguities:
$\Lform\to\Lform+\dd\bm{\mu}$ and $\tht\to\tht+\delta\bm{\mu}+\dd\bm{\nu}$.
These induce shifts in $(\J_\xi,\Q_\xi)$ and in the definition of the Hamiltonian by boundary terms.
The choice of $\Bnd$ (including the GHY term and a reference subtraction/counterterm) fixes both integrability and the zero-point of quasi-local energy.
\section{Extrinsic curvature on \texorpdfstring{$S^3$}{S3}}
For $\chi=\mathrm{const}$ spheres, the unit normal within $\Sigma$ is
\begin{equation}
n_i=\frac{1}{a}\,(\partial/\partial\chi)_i.
\end{equation}
A direct computation yields
\begin{equation}
k=\sigma^{AB}\nabla_A n_B=\frac{2}{a}\cot\chi,
\end{equation}
where $\sigma_{AB}$ is the induced metric on the 2-sphere $\partial D$ and $\nabla$ is the Levi--Civita connection of $\Sigma$.
\section{\texorpdfstring{$\zeta_{\goldconst}$}{zeta_phi_g} spectral overlay and modular reflection map (model-dependent)}
This appendix records a phenomenological overlay used in $\goldconst$-hierarchical modeling; it does not modify any GR statements in the main text and is logically independent of the quasi-local derivations.
Define an energy ladder
\begin{equation}
E^{(\zeta)}_n := \frac{\hbar c}{\ell_p\,\goldconst^n},
\end{equation}
with corresponding $m^{(\zeta)}_n=E^{(\zeta)}_n/c^2$.
Fix a distinguished index $N$ and define a reflection map $n':=N-n$.
\section{Explicit Killing generators on \texorpdfstring{$S^3$}{S3} (embedding form)}
Let $S^3$ be the unit sphere embedded in $\mathbb{R}^4$ with Cartesian coordinates $X^A$ obeying $X^AX^A=1$.
The full isometry group is $\mathrm{SO}(4)$, generated by six Killing fields $J_{AB}$:
\begin{equation}
J_{AB} := X^A\frac{\partial}{\partial X^B}-X^B\frac{\partial}{\partial X^A},
\qquad 1\le A<B\le 4.
\end{equation}
%------------------ Bibliography ------------------
\begin{thebibliography}{99}
\bibitem{Noether1918}
E.~Noether,
\newblock Invariante Variationsprobleme,
\newblock Nachr.\ Ges.\ Wiss.\ G\"ottingen, Math.-Phys.\ Kl.\ \textbf{1918}, 235 (1918).
\bibitem{Wald1993}
R.~M.~Wald,
\newblock Black hole entropy is the Noether charge,
\newblock Phys.\ Rev.\ D \textbf{48}, R3427 (1993).
\bibitem{IyerWald1994}
V.~Iyer and R.~M.~Wald,
\newblock Some properties of Noether charge and a proposal for dynamical black hole entropy,
\newblock Phys.\ Rev.\ D \textbf{50}, 846 (1994).
\bibitem{BrownYork1993}
J.~D.~Brown and J.~W.~York,
\newblock Quasilocal energy and conserved charges derived from the gravitational action,
\newblock Phys.\ Rev.\ D \textbf{47}, 1407 (1993).
\bibitem{MisnerSharp1964}
C.~W.~Misner and D.~H.~Sharp,
\newblock Relativistic equations for adiabatic, spherically symmetric gravitational collapse,
\newblock Phys.\ Rev.\ \textbf{136}, B571 (1964).
\bibitem{Szabados2009}
L.~B.~Szabados,
\newblock Quasi-local energy-momentum and angular momentum in general relativity,
\newblock Living Rev.\ Relativ.\ \textbf{12}, 4 (2009).
\bibitem{RoserValentini2014}
P.~Roser and A.~Valentini,
\newblock Classical and quantum cosmology with York time,
\newblock arXiv:1406.2036 [gr-qc] (2014).
\bibitem{Afshar2009}
M.~M.~Afshar,
\newblock Quasilocal energy in FRW cosmology,
\newblock arXiv:0903.3982 [gr-qc] (2009).
\bibitem{BrownKuchar1995}
J.~D.~Brown and K.~V.~Kucha\v{r},
\newblock Dust as a standard of space and time in canonical quantum gravity,
\newblock Phys.\ Rev.\ D \textbf{51}, 5600 (1995).
\bibitem{Dittrich2006}
B.~Dittrich,
\newblock Partial and complete observables for Hamiltonian constrained systems,
\newblock Class.\ Quant.\ Grav.\ \textbf{23}, 6155 (2006).


\bibitem{LeeWald1990}
J.~Lee and R.~M.~Wald,
\newblock Local symmetries and constraints,
\newblock J.\ Math.\ Phys.\ \textbf{31}, 725 (1990).

\bibitem{WaldZoupas2000}
R.~M.~Wald and A.~Zoupas,
\newblock A general definition of ``conserved quantities'' in general relativity and other theories of gravity,
\newblock Phys.\ Rev.\ D \textbf{61}, 084027 (2000).

\bibitem{GibbonsHawking1977}
G.~W.~Gibbons and S.~W.~Hawking,
\newblock Action integrals and partition functions in quantum gravity,
\newblock Phys.\ Rev.\ D \textbf{15}, 2752 (1977).

\bibitem{York1972}
J.~W.~York,
\newblock Role of conformal three-geometry in the dynamics of gravitation,
\newblock Phys.\ Rev.\ Lett.\ \textbf{28}, 1082 (1972).

\bibitem{Kodama1980}
H.~Kodama,
\newblock Conserved energy flux for the spherically symmetric system and the backreaction problem in the black hole evaporation,
\newblock Prog.\ Theor.\ Phys.\ \textbf{63}, 1217 (1980).

\bibitem{Hayward1996}
S.~A.~Hayward,
\newblock Gravitational energy in spherical symmetry,
\newblock Phys.\ Rev.\ D \textbf{53}, 1938 (1996).

\end{thebibliography}
\end{document}
